\documentclass[12pt]{article}
\usepackage[utf8]{inputenc}
\usepackage[T1]{fontenc}
\usepackage[french]{babel}
\usepackage{amsmath, amssymb}
\usepackage{geometry}
\geometry{a4paper, margin=2.5cm}

\title{Formalisation et Dualité du Problème de Covoiturage Dynamique}
\author{Modèle de Maximisation de Valeur Économique}
\date{\today}

\begin{document}

\maketitle

\section{Introduction}
Ce document présente la formalisation mathématique d'un problème de covoiturage dynamique visant à maximiser la valeur économique totale générée par le système, ainsi que sa formulation duale pour l'analyse des prix d'ombre.

\section{Paramètres et Variables}

\subsection{Ensembles et Indices}
\begin{itemize}
    \item $T$ : Ensemble des instants de temps discrets $t$.
    \item $P$ : Ensemble des passagers $p$.
    \item $D'$ : Ensemble étendu des conducteurs $d$ (incluant les véhicules fantômes).
\end{itemize}

\subsection{Paramètres}
\begin{itemize}
    \item $\Phi(t, p, d)$ : Valeur économique créée par le matching du passager $p$ avec le conducteur $d$ à l'instant $t$.
    \item $q_d$ : Capacité totale du véhicule $d$ (incluant le conducteur).
    \item $t_p^h$ : Heure de départ (disponibilité) du passager $p$.
    \item $\mathcal{A}$ : Ensemble des triplets $(t, p, d)$ admissibles satisfaisant les contraintes de compatibilité sociale ($penalty = 0$).
\end{itemize}

\subsection{Variable de décision Primal}
\begin{itemize}
    \item $x_{t,p,d} \in \{0, 1\}$ : Vaut 1 si le passager $p$ est transporté par le conducteur $d$ à l'instant $t$, 0 sinon.
\end{itemize}

\section{Problème Primal (Maximisation Sociale)}

L'objectif est de maximiser la valeur totale sur toute la durée de la simulation sous les contraintes de capacité, d'unicité et de disponibilité temporelle.

\begin{equation*}
\begin{aligned}
& \underset{x}{\text{Maximiser}} && Z_P = \sum_{(t,p,d) \in \mathcal{A}} \Phi(t, p, d) \cdot x_{t,p,d} \\
& \text{sous les contraintes :} && \\
& (1) \quad \sum_{p \in P} x_{t,p,d} \leq q_d - 1 && \forall t \in T, \forall d \in D' \quad [\lambda_{t,d}] \\
& (2) \quad \sum_{d \in D'} x_{t,p,d} \leq 1 && \forall t \in T, \forall p \in P \quad [\pi_{t,p}] \\
& (3) \quad x_{t,p,d} = 0 && \forall (t,p,d) \text{ tels que } t < t_p^h \\
& (4) \quad x_{t,p,d} \geq 0 &&
\end{aligned}
\end{equation*}

\section{Problème Dual (Prix de Marché)}

Le dual cherche à minimiser la valeur totale des ressources du système (sièges et disponibilité des passagers).

\begin{equation*}
\begin{aligned}
& \underset{\lambda, \pi}{\text{Minimiser}} && Z_D = \sum_{t \in T} \left( \sum_{d \in D'} (q_d - 1) \lambda_{t,d} + \sum_{p \in P} \pi_{t,p} \right) \\
& \text{sous les contraintes :} && \\
& \forall (t, p, d) \in \mathcal{A} \text{ avec } t \geq t_p^h : && \\
& \lambda_{t,d} + \pi_{t,p} \geq \Phi(t, p, d) \\
& \lambda_{t,d} \geq 0, \quad \pi_{t,p} \geq 0
\end{aligned}
\end{equation*}

\section{Interprétation Économique}

\begin{itemize}
    \item \textbf{Prix d'ombre des places ($\lambda_{t,d}$)} : Représente la valeur de rareté d'un siège dans le véhicule $d$ à l'instant $t$. Si le véhicule n'est pas saturé, $\lambda_{t,d} = 0$.
    \item \textbf{Rente du passager ($\pi_{t,p}$)} : Représente la valeur intrinsèque du passager pour le système à l'instant $t$, une fois déduite la "location" de sa place.
    \item \textbf{Complémentarité} : À l'optimum, un passager est matché ($x_{t,p,d} = 1$) uniquement si la somme des valeurs de ses ressources ($\lambda + \pi$) est exactement égale à la valeur $\Phi$ qu'il génère.
\end{itemize}

\end{document}
